\section{Limitations}\label{sec:limitations}
CRR is a reliability framework, not a universal accuracy replacement or an MVTec AUROC state-of-the-art claim. Official SubspaceAD remains a strong ranking baseline; our contribution concerns compact calibration, operational diagnostics, and source-assisted routing.

CRR is not adversarially robust. Label-aware PCA-residual FGSM stress tests cause severe degradation and depend on the attack objective, so corruption robustness must not be extrapolated to adversarial robustness.

Conformal validity depends on representative normal support data. LOIO residuals and full-support test residuals are asymmetric, and the normal p-values lie on a finite grid. Section~\ref{sec:uniformity} therefore reports clustered, discrete-reference diagnostics rather than treating a pooled iid uniformity test as exact. Without the relevant exchangeability conditions, target false-alarm control is not guaranteed.

The reported SC3R tables are historical empirical evidence. Candidate selection and source assessment were not independent, and their bootstrap intervals were pointwise rather than simultaneous. The $k{=}8$ false-alarm budget passed only at $\alpha{=}0.05$; at $\alpha{=}0.10$ the pooled rate was 0.121 against a 0.120 budget with 74\% no-harm, concentrated in Gaussian noise and JPEG. VisA no-harm fell to 78\% at $\alpha{=}0.20$, and the single MVTec-to-VisA direction traded power for lower observed false alarms. Source certification alone does not imply unconditional target-domain control. SC3R also requires multiple source categories, while matched-condition routing assumes condition metadata or a separately evaluated condition detector.

Finally, ECE for conformal-derived scores is prevalence-sensitive. Calibration numbers on anomaly-balanced benchmarks should not be interpreted as deployment reliability; false-alarm rate, power, precision, and risk--coverage are the primary operational summaries.
